%% ----------------------------------------------------------------------------
%%  §8.2  Somme integrali e integrale definito
%% ----------------------------------------------------------------------------
\section{Somme integrali e integrale definito}
\label{sec:integrale-definito}

Generalizziamo il procedimento del paragrafo precedente a una funzione qualsiasi,
purché \emph{limitata}. In tutto il paragrafo $[a,b]$ è un intervallo chiuso e limitato,
con $a<b$, e $f:[a,b]\to\R$ è una funzione limitata: esistono due costanti $m$, $M$ tali
che
\begin{equation}\label{eq:f-limitata}
  m\le f(x)\le M \qquad\perogni x\in[a,b] .
\end{equation}
Rispetto al settore di parabola cambiano due cose. La suddivisione dell'intervallo non
sarà più necessariamente in parti uguali; e, poiché $f$ non è più monotòna, l'altezza
dei rettangoli inscritti e circoscritti non sarà più il valore di $f$ in un estremo,
ma l'estremo inferiore e l'estremo superiore di $f$ nell'intervallo.

\subsection{Partizioni e somme integrali}

\begin{definizione}[Partizione]\label{def:partizione}\index{partizione}
Si chiama \emph{partizione} dell'intervallo $[a,b]$ un insieme finito
$P=\set{x_{0},x_{1},\dots,x_{n}}$ di punti, con $n\in\N$, tali che
\begin{equation}\label{eq:partizione}
  a=x_{0}<x_{1}<\dots<x_{k-1}<x_{k}<\dots<x_{n}=b .
\end{equation}
Gli $n+1$ punti individuano $n$ intervalli $I_{k}=[x_{k-1},x_{k}]$, $k=1,\dots,n$, di
ampiezza $\Delta x_{k}=x_{k}-x_{k-1}$. Il numero $\abs P=\max_{k}\Delta x_{k}$ si chiama
\emph{ampiezza} della partizione. Se $P$ e $Q$ sono partizioni di $[a,b]$ e
$P\subseteq Q$, si dice che $Q$ è \emph{più fine} di $P$, o che è un
\emph{raffinamento} di $P$.
\end{definizione}

Le ampiezze degli intervalli hanno per somma la lunghezza dell'intervallo:
$\sum_{k=1}^{n}\Delta x_{k}=x_{n}-x_{0}=b-a$ (ancora una somma telescopica). L'esempio
più semplice è la \emph{partizione uniforme} in $n$ parti uguali, con
$x_{k}=a+k\frac{b-a}n$ e $\abs P=\frac{b-a}n$, che abbiamo usato nel paragrafo
precedente. Se $P$ e $Q$ sono due partizioni di $[a,b]$, l'insieme
\begin{equation}\label{eq:unione-partizioni}
  R=P\cup Q ,
\end{equation}
ottenuto prendendo insieme i punti di $P$ e quelli di $Q$ (e ordinandoli), è ancora una
partizione di $[a,b]$, più fine sia di $P$ sia di $Q$.

\begin{definizione}[Somme integrali]\label{def:somme-integrali}\index{somma integrale}
Sia $P=\set{x_{0},\dots,x_{n}}$ una partizione di $[a,b]$. Per $k=1,\dots,n$ poniamo
\begin{equation}\label{eq:mk-Mk}
\begin{aligned}
  m_{k}&=\inf\set{f(x) \tc x\in[x_{k-1},x_{k}]},\\
  M_{k}&=\sup\set{f(x) \tc x\in[x_{k-1},x_{k}]} .
\end{aligned}
\end{equation}
Si chiamano \emph{somma integrale inferiore} e \emph{somma integrale superiore} di $f$
relative a $P$ i numeri
\begin{equation}\label{eq:somme-integrali}
  s(P)=\sum_{k=1}^{n}m_{k}\,(x_{k}-x_{k-1}), \qquad
  S(P)=\sum_{k=1}^{n}M_{k}\,(x_{k}-x_{k-1}) .
\end{equation}
Quando occorre ricordare la funzione si scrive $s(P,f)$ e $S(P,f)$.
\end{definizione}

Gli estremi $m_{k}$ e $M_{k}$ esistono finiti per il \cref{teo:esistenza-sup}, perché
l'insieme dei valori di $f$ in $[x_{k-1},x_{k}]$ è non vuoto e limitato; inoltre, per la
\eqref{eq:f-limitata}, $m\le m_{k}\le M_{k}\le M$. Moltiplicando per $\Delta x_{k}>0$ e
sommando si ottiene la stima
\begin{equation}\label{eq:stime-somme}
  m(b-a)\le s(P)\le S(P)\le M(b-a) \qquad\text{per ogni partizione } P .
\end{equation}
Se $f\ge0$, le somme integrali hanno il chiaro significato geometrico di somme delle
aree dei rettangoli rispettivamente inscritti nella regione sotto il grafico e
circoscritti a essa (\cref{fig:somme-integrali}). Si noti però che $s(P)$ e $S(P)$ sono
definite indipendentemente da ogni significato geometrico, anche se $f$ cambia segno.
Per $f(x)=x^{2}$ in $[0,b]$ e la partizione uniforme si ha $m_{k}=x_{k-1}^{2}$ e
$M_{k}=x_{k}^{2}$, e ritroviamo le somme $s_{n}$ e $S_{n}$ del \cref{sec:esaustione}.

\begin{figure}[htb]
\centering
\begin{tikzpicture}
\begin{axis}[width=0.9\textwidth, height=5.6cm, axis lines=left, xmin=0, xmax=4.95,
  ymin=0, ymax=3.1, xtick={0.3,1.3,2.5,3.3,4.5},
  xticklabels={$a=x_{0}$,$x_{1}$,$x_{2}$,$x_{3}$,$b=x_{4}$},
  ytick={1.345,2.424,2.733}, yticklabels={$m_{1}$,$m_{2}=M_{1}$,$M_{2}$},
  tick label style={font=\scriptsize}, clip=false]
  \fill[blu!20] (axis cs:0.3,0) rectangle (axis cs:1.3,1.345);
  \fill[blu!20] (axis cs:1.3,0) rectangle (axis cs:2.5,2.424);
  \fill[blu!20] (axis cs:2.5,0) rectangle (axis cs:3.3,1.899);
  \fill[blu!20] (axis cs:3.3,0) rectangle (axis cs:4.5,1.413);
  \draw[rame, thick] (axis cs:0.3,0) rectangle (axis cs:1.3,2.424);
  \draw[rame, thick] (axis cs:1.3,0) rectangle (axis cs:2.5,2.733);
  \draw[rame, thick] (axis cs:2.5,0) rectangle (axis cs:3.3,2.569);
  \draw[rame, thick] (axis cs:3.3,0) rectangle (axis cs:4.5,1.899);
  \addplot[curva, domain=0.2:4.65, samples=150] {1.7+0.8*sin(deg(1.3*x-0.9))+0.12*x};
  \addplot[tratteggio] coordinates {(0,1.345) (0.3,1.345)};
  \addplot[tratteggio] coordinates {(0,2.424) (1.3,2.424)};
  \addplot[tratteggio] coordinates {(0,2.733) (1.3,2.733)};
  \node[font=\scriptsize, blu, anchor=west] at (axis cs:4.67,1.53) {$y=f(x)$};
\end{axis}
\end{tikzpicture}
\caption{Somme integrali relative a una partizione non uniforme con $n=4$ intervalli.
L'area azzurra è la somma inferiore $s(P)$, quella racchiusa dai contorni in color rame
è la somma superiore $S(P)$. Nel secondo intervallo il massimo $M_{2}$ è assunto in un
punto interno, nel quarto il minimo $m_{4}$: le altezze dei rettangoli non sono, in
generale, i valori di $f$ negli estremi.}
\label{fig:somme-integrali}
\end{figure}

Il risultato seguente afferma che raffinando una partizione la somma inferiore
\emph{aumenta} (o resta uguale) e quella superiore \emph{diminuisce} (o resta uguale): le
due stime migliorano. Ne segue che \emph{ogni} somma inferiore è minore o uguale di
\emph{ogni} somma superiore, anche relativa a un'altra partizione.

\begin{lemma}[Raffinamento delle partizioni]\label{lem:raffinamento}
Siano $P$ e $Q$ partizioni di $[a,b]$.
\begin{enumerate}[label=\textup{(\alph*)}]
  \item Se $P\subseteq Q$, allora $s(P)\le s(Q)\le S(Q)\le S(P)$.
  \item In ogni caso, posto $R=P\cup Q$, risulta
  \begin{equation}\label{eq:s-minore-S}
    s(P)\le s(R)\le S(R)\le S(Q) .
  \end{equation}
\end{enumerate}
\end{lemma}

\begin{schema}
\Ipotesi{$f$ limitata in $[a,b]$; \ $P$, $Q$ partizioni di $[a,b]$.}
\Tesi{(a) $P\subseteq Q \implica s(P)\le s(Q)\le S(Q)\le S(P)$; \
(b) $s(P)\le S(Q)$ per ogni coppia $P$, $Q$.}
\end{schema}

\begin{strategia}
Il caso essenziale è quello in cui $Q$ ha un solo punto $\bar x$ in più di $P$. Allora
le due somme inferiori differiscono soltanto per i termini relativi all'intervallo di
$P$ che contiene $\bar x$, e tutto si riduce a un'osservazione sugli estremi inferiori:
l'estremo inferiore di $f$ su un insieme più piccolo è \emph{maggiore o uguale} di quello
sull'insieme più grande. Il caso generale si ottiene aggiungendo i punti di $Q$ uno alla
volta. La (b) segue dalla (a) applicata a $R$, che raffina sia $P$ sia $Q$.
\end{strategia}

\begin{dimostrazione}
(a) Supponiamo dapprima che $Q=P\cup\set{\bar x}$, con $x_{k-1}<\bar x<x_{k}$ per un certo
$k$, e poniamo
\[
  m'=\inf\set{f(x) \tc x\in[x_{k-1},\bar x]}, \qquad
  m''=\inf\set{f(x) \tc x\in[\bar x,x_{k}]} .
\]
Gli insiemi dei valori di $f$ in $[x_{k-1},\bar x]$ e in $[\bar x,x_{k}]$ sono contenuti
nell'insieme dei valori di $f$ in $[x_{k-1},x_{k}]$; poiché $m_{k}$ è un minorante di
quest'ultimo, è un minorante anche dei due sottoinsiemi, e quindi è minore o uguale dei
loro estremi inferiori: $m'\ge m_{k}$, $m''\ge m_{k}$. Le somme $s(Q)$ e $s(P)$ hanno in
comune tutti i termini tranne quelli relativi a $[x_{k-1},x_{k}]$, e dunque
\[
\begin{aligned}
  s(Q)-s(P)&=m'(\bar x-x_{k-1})+m''(x_{k}-\bar x)-m_{k}(x_{k}-x_{k-1})\\
  &\ge m_{k}\bigl[(\bar x-x_{k-1})+(x_{k}-\bar x)-(x_{k}-x_{k-1})\bigr]=0 .
\end{aligned}
\]
In modo analogo, gli estremi superiori sui due sottointervalli sono minori o uguali di
$M_{k}$, e si ottiene $S(Q)\le S(P)$. Se $Q$ contiene $p\ge2$ punti che non stanno in $P$,
si costruisce una catena $P=Q_{0}\subset Q_{1}\subset\dots\subset Q_{p}=Q$ in cui ciascuna
partizione si ottiene dalla precedente aggiungendo un solo punto; applicando il caso già
dimostrato a ogni anello della catena (cioè ragionando per induzione su $p$) si ottiene
$s(P)\le s(Q)$ e $S(Q)\le S(P)$. Infine $s(Q)\le S(Q)$ perché $m_{k}\le M_{k}$ per ogni
intervallo di $Q$.

(b) La partizione $R=P\cup Q$ è più fine sia di $P$ sia di $Q$. Applicando la (a) alla
coppia $P\subseteq R$ si ha $s(P)\le s(R)\le S(R)$, e applicandola alla coppia
$Q\subseteq R$ si ha $S(R)\le S(Q)$.
\end{dimostrazione}

\begin{figure}[htb]
\centering
\begin{tikzpicture}
\begin{axis}[width=0.62\textwidth, height=5cm, axis lines=left, xmin=0.4, xmax=4.1,
  ymin=0, ymax=4.3, xtick={1,2.4,3.6}, xticklabels={$x_{k-1}$,$\bar x$,$x_{k}$},
  ytick={0.85,2.04,3.84}, yticklabels={$m_{k}$,$m''=M'$,$M_{k}$},
  tick label style={font=\scriptsize}, clip=false]
  \fill[blu!20] (axis cs:1,0) rectangle (axis cs:3.6,0.85);
  \fill[verde!35] (axis cs:2.4,0.85) rectangle (axis cs:3.6,2.04);
  \fill[pattern=north east lines, pattern color=rame] (axis cs:1,2.04) rectangle
    (axis cs:2.4,3.84);
  \draw[rame, thick] (axis cs:1,0) rectangle (axis cs:3.6,3.84);
  \draw[grigio] (axis cs:2.4,0) -- (axis cs:2.4,3.84);
  \addplot[curva, domain=0.8:3.75, samples=80] {0.6+0.25*x^2};
  \addplot[tratteggio] coordinates {(0.4,2.04) (2.4,2.04)};
  \node[font=\scriptsize, verde!70!black] at (axis cs:3.0,1.45) {$+$};
  \node[font=\scriptsize, rame, fill=white, inner sep=1pt] at (axis cs:1.7,2.94) {$-$};
\end{axis}
\end{tikzpicture}
\caption{Aggiungendo il punto $\bar x$ all'intervallo $[x_{k-1},x_{k}]$, la somma
inferiore guadagna il rettangolo verde, perché in $[\bar x,x_{k}]$ l'estremo inferiore
sale da $m_{k}$ a $m''$; la somma superiore perde il rettangolo tratteggiato, perché in
$[x_{k-1},\bar x]$ l'estremo superiore scende da $M_{k}$ a $f(\bar x)=m''$.}
\label{fig:raffinamento}
\end{figure}

\subsection{L'integrale definito}

Indichiamo con $A$ l'insieme numerico costituito dalle somme inferiori $s(P)$, al variare
di $P$ tra tutte le partizioni di $[a,b]$, e con $B$ l'insieme delle corrispondenti somme
superiori:
\begin{equation}\label{eq:insiemi-A-B}
\begin{aligned}
  A&=\set{s(P) \tc P \text{ partizione di }[a,b]},\\
  B&=\set{S(P) \tc P \text{ partizione di }[a,b]} .
\end{aligned}
\end{equation}
Per il \cref{lem:raffinamento}(b), $a'\le b'$ per ogni $a'\in A$ e ogni $b'\in B$: i due
insiemi sono \emph{separati}. Per l'assioma di completezza (\cref{ass:completezza})
esiste almeno un numero reale $c$ maggiore o uguale di tutti gli elementi di $A$ e
minore o uguale di tutti gli elementi di $B$, cioè un \emph{elemento separatore}:
\[
  s(P)\le c\le S(Q) \qquad\text{per ogni coppia di partizioni } P, Q .
\]
In generale l'elemento separatore non è unico. Per descrivere tutti gli elementi
separatori introduciamo i numeri
\begin{equation}\label{eq:integrali-inf-sup}
\begin{aligned}
  s(f)&=\sup\set{s(P) \tc P \text{ partizione di }[a,b]},\\
  S(f)&=\inf\set{S(P) \tc P \text{ partizione di }[a,b]},
\end{aligned}
\end{equation}
detti rispettivamente \emph{integrale inferiore} e \emph{integrale superiore} di $f$.
Esistono finiti perché $A$ è non vuoto e limitato superiormente (da ogni elemento di
$B$), e $B$ è non vuoto e limitato inferiormente (da ogni elemento di $A$).

\begin{proposizione}[Gli elementi separatori]\label{prop:separatori}
Risulta $s(f)\le S(f)$, e un numero $c$ è elemento separatore tra $A$ e $B$ se e solo se
$s(f)\le c\le S(f)$.
\end{proposizione}

\begin{schema}
\Ipotesi{$f$ limitata in $[a,b]$; \ $A$, $B$ come in \eqref{eq:insiemi-A-B}.}
\Tesi{$s(f)\le S(f)$; \ gli elementi separatori sono tutti e soli i numeri
dell'intervallo $[s(f),S(f)]$.}
\end{schema}

\begin{strategia}
Tutto discende dalla definizione di estremo superiore come \emph{minimo dei maggioranti}
e di estremo inferiore come \emph{massimo dei minoranti}: un elemento separatore è un
maggiorante di $A$ che è anche un minorante di $B$.
\end{strategia}

\begin{dimostrazione}
Ogni elemento $S(Q)$ di $B$ è un maggiorante di $A$, e quindi è maggiore o uguale del
minimo dei maggioranti: $s(f)\le S(Q)$ per ogni $Q$. Dunque $s(f)$ è un minorante di $B$,
e quindi è minore o uguale del massimo dei minoranti: $s(f)\le S(f)$. Inoltre $c$ è
elemento separatore se e solo se è un maggiorante di $A$ e un minorante di $B$, cioè se
e solo se $c\ge\sup A=s(f)$ e $c\le\inf B=S(f)$.
\end{dimostrazione}

\begin{figure}[htb]
\centering
\begin{tikzpicture}[scale=1]
  \draw[-Latex, grigio] (0,1.1) -- (9.4,1.1);
  \foreach \x in {0.5,1.4,2.2,2.9,3.4,3.75,3.95} \fill[blu] (\x,1.1) circle (1.6pt);
  \foreach \x in {8.7,7.6,6.5,5.6,5.0,4.6,4.35} \fill[rame] (\x,1.1) circle (1.6pt);
  \draw[verde, very thick] (4.15,0.95) -- (4.15,1.25);
  \node[font=\scriptsize, above] at (4.15,1.25) {$s(f)=S(f)$};
  \node[font=\scriptsize, blu, below] at (2.0,1.0) {somme inferiori $s(P)$};
  \node[font=\scriptsize, rame, below] at (6.6,1.0) {somme superiori $S(P)$};
  \node[font=\scriptsize, anchor=west] at (9.5,1.1) {$f$ integrabile};
  \draw[-Latex, grigio] (0,0) -- (9.4,0);
  \foreach \x in {0.5,1.3,1.9,2.4,2.75,2.95,3.05} \fill[blu] (\x,0) circle (1.6pt);
  \foreach \x in {8.7,7.9,7.0,6.3,5.85,5.6,5.5} \fill[rame] (\x,0) circle (1.6pt);
  \draw[verde, very thick] (3.2,-0.15) -- (3.2,0.15);
  \draw[verde, very thick] (5.35,-0.15) -- (5.35,0.15);
  \fill[verde!25] (3.2,-0.06) rectangle (5.35,0.06);
  \node[font=\scriptsize, below] at (3.2,-0.15) {$s(f)$};
  \node[font=\scriptsize, below] at (5.35,-0.15) {$S(f)$};
  \node[font=\scriptsize, anchor=west] at (9.5,0) {$f$ non integrabile};
\end{tikzpicture}
\caption{In alto, le somme inferiori e quelle superiori si addensano attorno a un unico
elemento separatore: $f$ è integrabile. In basso, tra le une e le altre resta un
intervallo $[s(f),S(f)]$ di elementi separatori: $f$ non è integrabile.}
\label{fig:separatori}
\end{figure}

\begin{definizione}[Integrale definito secondo Riemann]\label{def:integrale}
\index{integrale definito}\index{funzione!integrabile}
Si dice che la funzione limitata $f$ è \emph{integrabile secondo Riemann} in $[a,b]$ se
esiste un \emph{unico} elemento separatore tra gli insiemi $A$ e $B$ delle somme
inferiori e superiori, cioè se $s(f)=S(f)$. Tale elemento si chiama \emph{integrale
definito} di $f$ in $[a,b]$ e si indica con il simbolo
\begin{equation}\label{eq:integrale}
  \int_{a}^{b}f(x)\,dx=s(f)=S(f) ,
\end{equation}
che si legge «integrale da $a$ a $b$ di $f(x)$ in $dx$».
\end{definizione}

Il termine «definito» distingue questo integrale, che è un numero, dagli
\emph{integrali indefiniti} del prossimo capitolo, nei quali gli estremi $a$, $b$ non
sono fissati. Vediamo subito due esempi opposti.

\begin{esempio}[Funzioni costanti]\label{es:integrale-costante}
Se $f(x)=c$ per ogni $x\in[a,b]$, in ogni intervallo di ogni partizione si ha
$m_{k}=M_{k}=c$, e quindi $s(P)=S(P)=c\sum\Delta x_{k}=c(b-a)$ per ogni $P$. Gli insiemi
$A$ e $B$ si riducono entrambi al numero $c(b-a)$, che è l'unico elemento separatore:
\begin{equation}\label{eq:integrale-costante}
  \int_{a}^{b}c\,dx=c\,(b-a) .
\end{equation}
\end{esempio}

\begin{esempio}[Una funzione limitata non integrabile]\label{es:dirichlet-non-integrabile}
Consideriamo in $[0,1]$ la funzione di Dirichlet $D$ dell'\cref{oss:dirichlet}, che
vale $1$ nei razionali e $0$ negli irrazionali. Per la densità dei razionali e degli
irrazionali (\cref{teo:densita-Q,cor:densita-irrazionali}), ogni intervallo
$[x_{k-1},x_{k}]$, con $x_{k-1}<x_{k}$, contiene sia punti razionali sia punti
irrazionali, e quindi $m_{k}=0$, $M_{k}=1$. Ne segue
\[
  s(P)=0, \qquad S(P)=\sum_{k=1}^{n}\Delta x_{k}=1 \qquad\text{per ogni partizione } P,
\]
dunque $s(D)=0<1=S(D)$: ogni numero di $[0,1]$ è un elemento separatore, e $D$ non è
integrabile secondo Riemann. La limitatezza, da sola, non basta all'integrabilità.
\end{esempio}

\subsection{Il criterio di integrabilità}

Verificare l'integrabilità calcolando $s(f)$ e $S(f)$, cioè un estremo superiore e un
estremo inferiore su \emph{tutte} le partizioni, è in genere impraticabile. Il teorema
seguente riduce il problema alla ricerca di \emph{una sola} partizione per ogni
tolleranza assegnata.

\begin{teorema}[Criterio di integrabilità]\label{teo:criterio-integrabilita}
\index{criterio!di integrabilità}
Una funzione $f$ limitata in $[a,b]$ è integrabile secondo Riemann in $[a,b]$ se e solo
se per ogni $\epsilon>0$ esiste una partizione $P$ di $[a,b]$ tale che
\begin{equation}\label{eq:criterio-integrabilita}
  S(P)-s(P)<\epsilon .
\end{equation}
\end{teorema}

\begin{schema}
\Ipotesi{$f$ limitata in $[a,b]$.}
\Tesi{$f$ integrabile $\sse$ $\perogni\epsilon>0\ \esiste P$ : $S(P)-s(P)<\epsilon$.}
\end{schema}

\begin{strategia}
\emph{Necessità}: se $s(f)=S(f)$, per le proprietà caratteristiche dell'estremo superiore
e dell'estremo inferiore esistono una somma inferiore e una somma superiore vicine a
meno di $\frac\epsilon2$ al valore comune, eventualmente relative a partizioni diverse
$P$ e $Q$; la partizione $R=P\cup Q$ le migliora entrambe (\cref{lem:raffinamento}).
\emph{Sufficienza}: $s(f)$ e $S(f)$ sono intrappolati tra $s(P)$ e $S(P)$, quindi distano
meno di $\epsilon$; un numero non negativo minore di ogni $\epsilon>0$ è nullo.
\end{strategia}

\begin{dimostrazione}
\emph{Necessità.} Sia $f$ integrabile, cioè $s(f)=S(f)$, e sia $\epsilon>0$. Per la
caratterizzazione dell'estremo superiore (\cref{prop:caratt-sup}) esiste una partizione
$P$ con $s(f)-\frac\epsilon2<s(P)$, e per quella dell'estremo inferiore
\eqref{eq:caratt-inf} esiste una partizione $Q$ con $S(Q)<S(f)+\frac\epsilon2$. Posto
$R=P\cup Q$, per la \eqref{eq:s-minore-S}
\[
  s(f)-\frac\epsilon2<s(P)\le s(R)\le S(R)\le S(Q)<S(f)+\frac\epsilon2 ,
\]
e quindi, essendo $s(f)=S(f)$,
\[
  S(R)-s(R)<\Bigl(S(f)+\frac\epsilon2\Bigr)-\Bigl(s(f)-\frac\epsilon2\Bigr)=\epsilon .
\]
La partizione cercata è $R$.

\emph{Sufficienza.} Per ogni partizione $P$ si ha $s(P)\le s(f)$, perché $s(f)$ è un
maggiorante delle somme inferiori, e $S(f)\le S(P)$, perché $S(f)$ è un minorante delle
somme superiori. Insieme alla \cref{prop:separatori} si ottiene
$s(P)\le s(f)\le S(f)\le S(P)$, e quindi
\[
  0\le S(f)-s(f)\le S(P)-s(P) .
\]
Se vale l'ipotesi, per ogni $\epsilon>0$ possiamo scegliere $P$ in modo che il secondo
membro sia minore di $\epsilon$: il numero $S(f)-s(f)$, che non dipende da $\epsilon$, è
non negativo e minore di ogni numero positivo, e dunque è nullo. Perciò $s(f)=S(f)$, cioè
$f$ è integrabile.
\end{dimostrazione}

Nella pratica il criterio si usa quasi sempre attraverso una successione di partizioni,
come abbiamo fatto per la parabola.

\begin{corollario}[Integrale come limite di somme]\label{cor:limite-somme}
Sia $f$ limitata in $[a,b]$ e sia $P_{n}$ una successione di partizioni di $[a,b]$ tale
che $S(P_{n})-s(P_{n})\to0$. Allora $f$ è integrabile e
\begin{equation}\label{eq:limite-somme}
  \int_{a}^{b}f(x)\,dx=\lim_{n\to+\infty}s(P_{n})=\lim_{n\to+\infty}S(P_{n}) .
\end{equation}
\end{corollario}

\begin{schema}
\Ipotesi{$f$ limitata in $[a,b]$; \ $S(P_{n})-s(P_{n})\to0$.}
\Tesi{$f$ integrabile; \ $s(P_{n})\to\int_{a}^{b}f$ e $S(P_{n})\to\int_{a}^{b}f$.}
\end{schema}

\begin{strategia}
L'integrabilità è il criterio appena dimostrato. Poi l'integrale, come ogni elemento
separatore, sta tra $s(P_{n})$ e $S(P_{n})$, e il teorema dei carabinieri fa il resto.
\end{strategia}

\begin{dimostrazione}
Dato $\epsilon>0$, per la definizione di limite esiste $n$ con
$S(P_{n})-s(P_{n})<\epsilon$, e per il \cref{teo:criterio-integrabilita} $f$ è
integrabile. Posto $I=\int_{a}^{b}f(x)\,dx$, si ha $s(P_{n})\le I\le S(P_{n})$, quindi
\[
  0\le I-s(P_{n})\le S(P_{n})-s(P_{n}), \qquad 0\le S(P_{n})-I\le S(P_{n})-s(P_{n}) ,
\]
e per il teorema dei carabinieri (\cref{teo:carabinieri}) $I-s(P_{n})\to0$ e
$S(P_{n})-I\to0$.
\end{dimostrazione}

\begin{esempio}[Gli integrali di $x$ e di $x^{2}$]\label{es:integrale-x-x2}
Sia $f(x)=x$ in $[0,b]$, e sia $P_{n}$ la partizione uniforme con $x_{k}=\frac{kb}n$.
Poiché $f$ è crescente, $m_{k}=x_{k-1}$ e $M_{k}=x_{k}$; con la somma dei primi
$n$ numeri naturali (\cref{prop:gauss}),
\[
\begin{aligned}
  S(P_{n})&=\frac bn\sum_{k=1}^{n}\frac{kb}n=\frac{b^{2}}{n^{2}}\cdot\frac{n(n+1)}2
  =\frac{b^{2}}2\cdot\frac{n+1}n,\\
  s(P_{n})&=\frac{b^{2}}2\cdot\frac{n-1}n ,
\end{aligned}
\]
e $S(P_{n})-s(P_{n})=\frac{b^{2}}n\to0$. Per il \cref{cor:limite-somme}
\begin{equation}\label{eq:integrale-x}
  \int_{0}^{b}x\,dx=\frac{b^{2}}2 ,
\end{equation}
che è l'area del triangolo di vertici $(0,0)$, $(b,0)$, $(b,b)$, come deve essere. Allo
stesso modo, le somme \eqref{eq:Sn-parabola}--\eqref{eq:sn-parabola} del settore di
parabola hanno differenza $\frac{b^{3}}n\to0$, e quindi
\begin{equation}\label{eq:integrale-x2}
  \int_{0}^{b}x^{2}\,dx=\frac{b^{3}}3 .
\end{equation}
\end{esempio}

Per una funzione non monotòna il calcolo di $m_{k}$ e $M_{k}$ richiede un po' più di
lavoro. Se $f$ è derivabile, però, gli strumenti del \cref{cap:applicazioni-derivate}
bastano: in ogni intervallo $[x_{k-1},x_{k}]$ la funzione, continua, ha massimo e minimo
(\cref{teo:weierstrass}), e se uno di essi è assunto in un punto interno, lì la derivata
si annulla (\cref{teo:fermat}). Dunque $m_{k}$ e $M_{k}$ vanno cercati tra i valori negli
estremi e quelli nei punti critici interni. Il programma seguente applica questa
osservazione a $f(x)=x^{3}-3x$ in $[-2,2]$, i cui punti critici sono $\pm1$.

\begin{python}[Somme integrali di una funzione non monotòna]
def darboux(f, a, b, n, critici=()):
    """Somme s(P), S(P) per la partizione uniforme in n parti:
    m_k e M_k si cercano tra gli estremi dell'intervallo e i
    punti critici interni (Weierstrass + Fermat)."""
    x = [a + k * (b - a) / n for k in range(n + 1)]
    s = S = 0.0
    for k in range(1, n + 1):
        cand = [x[k - 1], x[k]]
        cand += [c for c in critici if x[k - 1] < c < x[k]]
        v = [f(t) for t in cand]
        s += min(v) * (x[k] - x[k - 1])
        S += max(v) * (x[k] - x[k - 1])
    return s, S

f = lambda x: x**3 - 3 * x            # punti critici: -1 e 1
for n in (3, 30, 300, 3000):
    s, S = darboux(f, -2, 2, n, critici=(-1, 1))
    print(f"n = {n:4d}: s = {s:9.5f}  S = {S:8.5f}"
          f"  S-s = {S - s:8.5f}  n(S-s) = {n * (S - s):5.2f}")
\end{python}
\begin{risultato}
n =    3: s =  -7.60494  S =  7.60494  S-s = 15.20988  n(S-s) = 45.63
n =   30: s =  -0.79826  S =  0.79826  S-s =  1.59652  n(S-s) = 47.90
n =  300: s =  -0.08000  S =  0.08000  S-s =  0.16000  n(S-s) = 48.00
n = 3000: s =  -0.00800  S =  0.00800  S-s =  0.01600  n(S-s) = 48.00
\end{risultato}

Le due somme si stringono attorno a $0$, il valore suggerito dalla simmetria ($f$ è
dispari, e le aree sopra e sotto l'asse si compensano), e la differenza decresce come
$1/n$. Il prodotto $n(S-s)$ si stabilizza intorno a $48$, e il motivo è semplice: $f$
sale da $-2$ a $2$ in $[-2,-1]$, scende a $-2$ in $[-1,1]$ e risale a $2$ in $[1,2]$;
in ogni intervallo $M_{k}-m_{k}$ è al più la variazione totale di $f$ in
quell'intervallo (la somma delle sue salite e delle sue discese), e la somma di tutte le
salite e discese in $[-2,2]$ è $4+4+4=12$. Quindi
$S-s=\sum(M_{k}-m_{k})\frac4n\le\frac{48}n$.

\subsection{Significato geometrico e notazioni}

\begin{definizione}[Rettangoloide]\label{def:rettangoloide}\index{rettangoloide}
Sia $f\ge0$ in $[a,b]$. L'insieme
\begin{equation}\label{eq:rettangoloide}
  R_{f}=\set{(x,y)\in\R^{2} \tc a\le x\le b,\ 0\le y\le f(x)}
\end{equation}
si chiama \emph{rettangoloide} (o \emph{trapezoide}) di base $[a,b]$ relativo a $f$.
\end{definizione}

Per ogni partizione $P$, la somma $s(P)$ è l'area di un plurirettangolo contenuto in
$R_{f}$, e $S(P)$ quella di un plurirettangolo che contiene $R_{f}$
(\cref{fig:somme-integrali}). Il criterio di integrabilità afferma che, se $f$ è
integrabile, si possono trovare un plurirettangolo contenuto in $R_{f}$ e uno che lo
contiene le cui aree differiscono per meno di un qualsiasi $\epsilon>0$ assegnato. È
dunque ragionevole \emph{definire} l'area del rettangoloide come l'unico numero che
separa le aree dei plurirettangoli inscritti da quelle dei plurirettangoli
circoscritti:
\begin{equation}\label{eq:area-rettangoloide}
  \text{area }R_{f}=\int_{a}^{b}f(x)\,dx \qquad (f\ge0 \text{ integrabile}) .
\end{equation}
La definizione è coerente con le aree elementari: per una funzione costante si ottiene
l'area del rettangolo, per $f(x)=x$ quella del triangolo, per $f(x)=x^{2}$ quella del
settore di parabola calcolata da Archimede. È il punto di partenza della teoria della
misura delle figure piane di Peano e Jordan.

Nella scrittura $\int_{a}^{b}f(x)\,dx$ i numeri $a$ e $b$ si chiamano \emph{estremi di
integrazione}, la funzione $f$ \emph{funzione integranda}, la variabile $x$
\emph{variabile di integrazione}. Il risultato è un numero e non dipende da $x$, che è
una variabile «muta», come l'indice di una sommatoria: $\int_{a}^{b}f(x)\,dx$,
$\int_{a}^{b}f(t)\,dt$ e $\int_{a}^{b}f(s)\,ds$ indicano lo stesso numero. È utile
considerare l'integrale anche quando il primo estremo non è minore del secondo: si pone
\begin{equation}\label{eq:estremi-invertiti}
  \int_{a}^{b}f(x)\,dx=-\int_{b}^{a}f(x)\,dx \quad (a>b), \qquad
  \int_{a}^{a}f(x)\,dx=0 .
\end{equation}

\begin{osservazione}[Somme di Riemann]\label{oss:somme-riemann}\index{somma di Riemann}
Data una partizione $P$, scegliamo in ciascun intervallo $[x_{k-1},x_{k}]$ un punto
$\xi_{k}$ arbitrario e formiamo la \emph{somma di Riemann}
\begin{equation}\label{eq:somma-riemann}
  \sigma(P,\xi)=\sum_{k=1}^{n}f(\xi_{k})\,(x_{k}-x_{k-1}) .
\end{equation}
Poiché $m_{k}\le f(\xi_{k})\le M_{k}$, si ha $s(P)\le\sigma(P,\xi)\le S(P)$; se $f$ è
integrabile, anche l'integrale sta tra $s(P)$ e $S(P)$, e quindi
\begin{equation}\label{eq:riemann-stima}
  \Bigl|\sigma(P,\xi)-\int_{a}^{b}f(x)\,dx\Bigr|\le S(P)-s(P) .
\end{equation}
Le somme di Riemann sono quelle che si calcolano in pratica, perché richiedono soltanto
valori di $f$ e non estremi inferiori e superiori: scegliendo come $\xi_{k}$ l'estremo
sinistro, l'estremo destro o il punto medio di ogni intervallo si ottengono le formule
dei rettangoli che studieremo nel \cref{sec:integrabilita-continue}. La notazione di
Leibniz ricorda proprio queste somme: il simbolo $\int$ è una «S» allungata, iniziale
di \emph{summa}, e $dx$ ricorda le ampiezze $\Delta x_{k}$.
\end{osservazione}

\begin{notastorica}[Leibniz, Cauchy, Riemann, Darboux]
Gottfried Wilhelm Leibniz introdusse il simbolo $\int$ in un manoscritto del 29 ottobre
1675, sostituendo l'abbreviazione \emph{omn.} (\emph{omnia}, la «somma di tutte» le
linee di una figura) con una S allungata, iniziale di \emph{summa}; fu poi Joseph
Fourier, intorno al 1820, a proporre di scrivere gli estremi in basso e in alto. Per i
successori di Leibniz, a partire da Johann Bernoulli ed Eulero, l'integrale divenne
soprattutto l'operazione inversa della derivazione. Il primo a definirlo in modo
autonomo, come limite di somme, fu Augustin-Louis Cauchy nel \emph{Résumé des leçons}
del 1823, per le funzioni continue e con le somme costruite sugli estremi sinistri: il
legame con la derivata, che per Leibniz era quasi una definizione, diventò così un
teorema, quello fondamentale del prossimo capitolo. Cauchy deduceva l'esistenza del
limite dal fatto che due somme relative a partizioni abbastanza fini differiscono poco,
un passaggio che richiede la completezza dei reali, allora non ancora formulata; inoltre
la somma dipende da tutta la partizione e non dalla sola ampiezza, e un limite «per
ampiezza tendente a zero» sfuggiva alle definizioni di limite note (per questo nel
\cref{teo:riemann-continue} useremo successioni di partizioni). Una teoria dei limiti
adatta a questi casi fu proposta solo nel 1922 da Eliakim H.~Moore e Herman L.~Smith.
Bernhard Riemann, nella dissertazione di abilitazione del 1854 sulle serie
trigonometriche (pubblicata postuma, a cura di Dedekind, nel 1867), si pose la domanda
dell'epigrafe e rispose con le somme \eqref{eq:somma-riemann} a punti $\xi_{k}$
arbitrari, ammettendo funzioni qualsiasi purché limitate; diede anche il criterio di
integrabilità. La formulazione con le somme inferiori e superiori, che abbiamo seguito,
è dovuta a Gaston Darboux (1875); le due definizioni sono equivalenti.
\end{notastorica}

\subsection{Approfondimento: l'integrale con le funzioni a scala}
\label{subsec:funzioni-scala}

Molti testi, tra cui \cite{giusti}, definiscono l'integrale in un modo che a prima vista
sembra diverso dal nostro: al posto delle somme inferiori e superiori usano le
\emph{funzioni a scala} che stanno sotto e sopra $f$. Mostriamo che le due definizioni
coincidono; la verifica richiede un po' di attenzione nei punti di suddivisione.

\begin{definizione}[Funzione a scala]\label{def:funzione-scala}\index{funzione!a scala}
Una funzione $\varphi:[a,b]\to\R$ si dice \emph{a scala} se esiste una partizione
$P=\set{x_{0},\dots,x_{n}}$ di $[a,b]$ tale che $\varphi$ è costante in ciascun
intervallo aperto $(x_{k-1},x_{k})$: $\varphi(x)=c_{k}$ per $x_{k-1}<x<x_{k}$. Una tale
partizione si dice \emph{adatta} a $\varphi$; i valori di $\varphi$ nei punti $x_{k}$
sono arbitrari. Il numero
\begin{equation}\label{eq:integrale-scala}
  I(\varphi)=\sum_{k=1}^{n}c_{k}\,(x_{k}-x_{k-1})
\end{equation}
si chiama \emph{integrale elementare} di $\varphi$.
\end{definizione}

Il numero $I(\varphi)$ non dipende dalla partizione adatta. Se si aggiunge a $P$ un
punto $\bar x\in(x_{k-1},x_{k})$, la partizione resta adatta e il termine
$c_{k}(x_{k}-x_{k-1})$ si spezza in $c_{k}(\bar x-x_{k-1})+c_{k}(x_{k}-\bar x)$, con la
stessa somma; aggiungendo i punti uno alla volta, come nel \cref{lem:raffinamento}, la
somma non cambia passando a un raffinamento. Se $P$ e $P'$ sono adatte, $P\cup P'$
raffina entrambe, e le tre somme coincidono. Se $f\ge0$ e $\varphi\le f$, $I(\varphi)$ è
l'area di un plurirettangolo contenuto nel rettangoloide di $f$, come una somma
inferiore.

Sia ora $f$ limitata in $[a,b]$, con $m\le f\le M$. Poniamo
\begin{equation}\label{eq:integrali-scala}
\begin{aligned}
  I_{-}(f)&=\sup\set{I(\psi) \tc \psi \text{ a scala},\ \psi\le f \text{ in }[a,b]},\\
  I^{+}(f)&=\inf\set{I(\varphi) \tc \varphi \text{ a scala},\ \varphi\ge f \text{ in }[a,b]} .
\end{aligned}
\end{equation}
Gli insiemi sono non vuoti, perché le costanti $m$ e $M$ sono funzioni a scala; e se
$\psi\le f$ è a scala, i suoi valori $c_{k}$ sono minori o uguali di $M$, quindi
$I(\psi)\le M(b-a)$. Analogamente $I(\varphi)\ge m(b-a)$, e i due numeri sono finiti.

\begin{proposizione}[Funzioni a scala e somme integrali]\label{prop:scala-darboux}
Se $f$ è limitata in $[a,b]$, allora $I_{-}(f)=s(f)$ e $I^{+}(f)=S(f)$. In particolare
$f$ è integrabile secondo Riemann se e solo se $I_{-}(f)=I^{+}(f)$, e in tal caso il
valore comune è $\int_{a}^{b}f(x)\,dx$.
\end{proposizione}

\begin{schema}
\Ipotesi{$f$ limitata in $[a,b]$, \ $m\le f\le M$.}
\Tesi{$I_{-}(f)=s(f)$ \ e \ $I^{+}(f)=S(f)$, \ con $I_{-}$, $I^{+}$ definiti dalla
\eqref{eq:integrali-scala}.}
\end{schema}

\begin{strategia}
Ogni somma inferiore è l'integrale elementare di una funzione a scala minorante:
quella che vale $m_{k}$ all'interno del $k$-esimo intervallo e coincide con $f$ nei punti
di suddivisione. Dunque $I_{-}(f)\ge s(f)$. Il viceversa è più delicato, perché una
funzione a scala minorante può superare $m_{k}$ all'interno di un intervallo, se $f$
assume valori piccoli soltanto negli estremi: si rimedia isolando i punti di
suddivisione in intervallini di ampiezza $\eta$, il cui contributo è proporzionale a
$\eta$. Per $I^{+}(f)$ basta passare a $-f$.
\end{strategia}

\begin{dimostrazione}
\emph{$I_{-}(f)\ge s(f)$.} Sia $P=\set{x_{0},\dots,x_{n}}$ una partizione, e sia
$\psi_{P}$ la funzione che vale $m_{k}$ in $(x_{k-1},x_{k})$ e $f(x_{k})$ nei punti
$x_{k}$. È una funzione a scala, con $P$ adatta, e $\psi_{P}\le f$: all'interno degli
intervalli perché $m_{k}$ è un minorante di $f$ in $[x_{k-1},x_{k}]$, nei punti di
suddivisione per costruzione. Dunque $s(P)=I(\psi_{P})\le I_{-}(f)$ per ogni $P$, e
passando all'estremo superiore $s(f)\le I_{-}(f)$.

\emph{$I_{-}(f)\le s(f)$.} Sia $\psi\le f$ a scala, con partizione adatta
$P=\set{x_{0},\dots,x_{n}}$ e valori $c_{k}$; come osservato, $c_{k}\le M$. Fissiamo
$\eta$ con $0<2\eta<\min_{k}\Delta x_{k}$ e sia $Q$ la partizione ottenuta aggiungendo a
$P$ i punti $x_{k-1}+\eta$ e $x_{k}-\eta$, per $k=1,\dots,n$. Ogni intervallo
$[x_{k-1},x_{k}]$ si divide così in due intervallini di ampiezza $\eta$ adiacenti agli
estremi e nell'intervallo centrale $J_{k}=[x_{k-1}+\eta,x_{k}-\eta]$. Poiché $J_{k}$ è
contenuto in $(x_{k-1},x_{k})$, dove $f\ge\psi=c_{k}$, l'estremo inferiore di $f$ in
$J_{k}$ è maggiore o uguale di $c_{k}$; negli intervallini è maggiore o uguale di $m$.
La partizione $Q$ è adatta a $\psi$, e con essa
\[
  I(\psi)=\sum_{k=1}^{n}\bigl(c_{k}\abs{J_{k}}+2c_{k}\eta\bigr), \qquad
  s(Q)\ge\sum_{k=1}^{n}\bigl(c_{k}\abs{J_{k}}+2m\eta\bigr),
\]
dove $\abs{J_{k}}$ è l'ampiezza di $J_{k}$. Sottraendo,
\[
  I(\psi)-s(Q)\le2\eta\sum_{k=1}^{n}(c_{k}-m)\le2n(M-m)\,\eta ,
\]
e quindi $I(\psi)\le s(f)+2n(M-m)\eta$. Qui $n$ dipende soltanto da $\psi$, mentre $\eta$
si può scegliere piccolo a piacere: il numero $I(\psi)-s(f)$ è minore o uguale di ogni
numero positivo, e quindi $I(\psi)\le s(f)$. Passando all'estremo superiore su $\psi$,
$I_{-}(f)\le s(f)$.

\emph{Integrali superiori.} Si ha $\varphi\ge f$ se e solo se $-\varphi\le-f$, e
$I(-\varphi)=-I(\varphi)$; quindi $I^{+}(f)=-I_{-}(-f)$. D'altra parte, in ogni
intervallo l'estremo inferiore di $-f$ è l'opposto dell'estremo superiore di $f$, per cui
$s(P,-f)=-S(P,f)$ e $s(-f)=-S(f)$. Applicando a $-f$ quanto già dimostrato,
$I^{+}(f)=-I_{-}(-f)=-s(-f)=S(f)$. L'ultima affermazione segue dalla
\cref{def:integrale}.
\end{dimostrazione}

\begin{osservazione}[Integrale delle funzioni a scala]\label{oss:integrale-scala}
Ogni funzione a scala $\varphi$ è integrabile, e $\int_{a}^{b}\varphi(x)\,dx=I(\varphi)$.
Infatti $\varphi$, che assume un numero finito di valori, è limitata, ed è minorante e
maggiorante di sé stessa; per la \eqref{eq:integrali-scala}, la
\cref{prop:scala-darboux} e la \cref{prop:separatori},
$I(\varphi)\le I_{-}(\varphi)=s(\varphi)\le S(\varphi)=I^{+}(\varphi)\le I(\varphi)$, e
tutte le disuguaglianze sono uguaglianze. Ad esempio, una funzione nulla tranne che in
un numero finito di punti è a scala, con integrale nullo: con la linearità del
\cref{sec:proprieta-integrale} si ritrova l'\cref{es:valore-in-un-punto}. Le due
definizioni sono dunque la stessa idea, l'\emph{unico elemento separatore}, vista con
due strumenti: le somme sulle partizioni o i plurirettangoli delle funzioni a scala.
\end{osservazione}

\begin{notastorica}[Il commerciante di Lebesgue]
Approssimare $f$ dal basso e dall'alto con funzioni costanti su intervalli è anche il
punto di partenza della teoria della misura. Giuseppe Peano nel 1887 e Camille Jordan
pochi anni dopo definirono la misura interna ed esterna di un insieme mediante i
plurirettangoli contenuti nell'insieme e quelli che lo contengono, come le somme
inferiori e superiori del rettangoloide; ma l'insieme dei razionali di $[0,1]$ ha misura
interna $0$ ed esterna $1$, e non è misurabile. Émile Borel, nel 1898, chiese alla misura
di essere additiva anche per successioni di insiemi disgiunti, e su questa base Henri
Lebesgue, nella tesi \emph{Intégrale, longueur, aire} (1902), costruì un integrale nuovo.
In una conferenza del 1926 ne spiegò l'idea con un'immagine: Riemann somma i valori di
$f$ nell'ordine in cui li incontra percorrendo l'intervallo, come un commerciante
disordinato che conta le monete nell'ordine in cui gli capitano in mano; Lebesgue
invece le raggruppa per valore, come il commerciante metodico che conta quante monete
ha di ciascun taglio. Si suddivide cioè l'insieme dei \emph{valori} di $f$, non
l'intervallo, e si approssima $f$ con funzioni costanti su insiemi più generali degli
intervalli, dei quali occorre conoscere la misura. Il guadagno non è soltanto una classe
più ampia di funzioni integrabili (la funzione di Dirichlet lo è, con integrale nullo),
ma soprattutto i teoremi di passaggio al limite sotto il segno di integrale, essenziali
per le serie di Fourier. Non tutti gli insiemi, però, sono misurabili: nel 1905
Giuseppe Vitali ne costruì uno che non lo è, servendosi dell'assioma della scelta,
formulato da Ernst Zermelo l'anno prima.
\end{notastorica}
